
\documentclass[UTF8,a4paper,11pt]{ctexart}

\usepackage{amsmath,amssymb,amsthm,bm}
\usepackage{geometry}
\usepackage{booktabs}
\usepackage{array}
\usepackage{physics}
\usepackage{mathtools}
\usepackage{algorithm}
\usepackage{algpseudocode}
\usepackage{xcolor}
\usepackage{hyperref}

\geometry{left=2.4cm,right=2.4cm,top=2.5cm,bottom=2.5cm}

\hypersetup{
  colorlinks=true,
  linkcolor=blue,
  citecolor=blue,
  urlcolor=blue
}

\newcommand{\eps}{\varepsilon}
\newcommand{\Roo}{R^{oo}}
\newcommand{\Rvv}{R^{vv}}
\newcommand{\Ria}{R^{IA}}
\newcommand{\Had}{\odot}
\newcommand{\trans}{\mathrm{T}}

\title{基于 RI 张量投影的四指标求和式高效矩阵--张量实现}
\author{}
\date{}

\begin{document}

\maketitle

\section{问题定义}

考虑如下线性变换：
\begin{equation}
Y_{ia}
=
\sum_{jbIA}
\left(
\sum_{\mu} R_{ij,\mu} R_{IA,\mu}
\right)
\left(
\sum_{\nu} R_{ab,\nu} R_{IA,\nu}
\right)
\left[
\frac{1}{
\omega - \eps_b + \eps_i
- \eps_A + \eps_I
}
-
\frac{1}{
\omega - \eps_a + \eps_j
- \eps_A + \eps_I
}
\right]
X_{jb}.
\label{eq:original}
\end{equation}

其中：
\begin{itemize}
  \item $i,j$ 是 occupied 轨道指标；
  \item $a,b$ 是 virtual 轨道指标；
  \item $I,A$ 是内部求和中的 occupied--virtual 指标；
  \item $\mu,\nu$ 是辅助基指标；
  \item $R$ 是三阶 RI 张量在不同二指标组合上的矩阵化表示；
  \item $X_{jb}$ 是输入向量，按 occupied--virtual pair 排列后可以视为矩阵；
  \item $Y_{ia}$ 是输出向量，也可以视为矩阵；
  \item $\omega$ 是外部频率参数。
\end{itemize}

实际目标是实现如下 matvec：
\begin{equation}
X_{jb} \longmapsto Y_{ia},
\end{equation}
并且避免显式形成四指标对象
\begin{equation}
\widetilde W_{ij,ab}
=
\sum_{IA}
\left(
\sum_{\mu} R_{ij,\mu} R_{IA,\mu}
\right)
\left(
\sum_{\nu} R_{ab,\nu} R_{IA,\nu}
\right)
\left[
\frac{1}{
\omega - \eps_b + \eps_i
- \eps_A + \eps_I
}
-
\frac{1}{
\omega - \eps_a + \eps_j
- \eps_A + \eps_I
}
\right].
\end{equation}

直接构造 $\widetilde W_{ij,ab}$ 的存储量为
\begin{equation}
O(n_o^2 n_v^2),
\end{equation}
其中 $n_o$ 是 occupied 轨道数，$n_v$ 是 virtual 轨道数。
这在实际 BSE/TDDFT/GW-BSE 类型程序中通常是不可接受的。
因此需要利用 RI 低秩结构，将原始四指标求和改写为一系列矩阵乘法、逐元素缩放和张量重排。

\section{合并指标与中间量定义}

将内部求和指标 $(I,A)$ 合并为一个指标
\begin{equation}
p=(I,A).
\end{equation}
定义
\begin{equation}
\Delta_p
=
\Delta_{IA}
=
\eps_I-\eps_A.
\end{equation}
于是式 \eqref{eq:original} 可以写为
\begin{equation}
Y_{ia}
=
\sum_{p}\sum_{jb}
U^p_{ij}
V^p_{ab}
\left[
\frac{1}{\omega+\eps_i-\eps_b+\Delta_p}
-
\frac{1}{\omega+\eps_j-\eps_a+\Delta_p}
\right]
X_{jb},
\label{eq:compact_original}
\end{equation}
其中定义两个 RI 投影中间量：
\begin{align}
U^p_{ij}
&=
\sum_{\mu} R_{ij,\mu}R_{p,\mu},
\label{eq:U_def}
\\
V^p_{ab}
&=
\sum_{\nu} R_{ab,\nu}R_{p,\nu}.
\label{eq:V_def}
\end{align}

这里 $R_{p,\mu}$ 即原来的 $R_{IA,\mu}$。

式 \eqref{eq:compact_original} 的关键结构是：
对于每一个固定的 $p=(I,A)$，四指标对象在 $(ij)$ 与 $(ab)$ 上是低秩外积形式：
\begin{equation}
U^p_{ij}V^p_{ab}.
\end{equation}
但是由于能量分母同时依赖外层指标，不能简单地把所有 $p$ 的贡献预先累加成一个静态矩阵。
正确做法是对每个 $p$ 或每个 $p$ 的 block 分别处理能量分母。

\section{第一项的矩阵化}

先考虑式 \eqref{eq:compact_original} 中括号内的第一项：
\begin{equation}
Y^{(1,p)}_{ia}
=
\sum_{jb}
U^p_{ij}V^p_{ab}
\frac{1}{\omega+\eps_i-\eps_b+\Delta_p}
X_{jb}.
\label{eq:first_term}
\end{equation}

定义分母矩阵
\begin{equation}
F^p_{ib}
=
\frac{1}{\omega+\eps_i-\eps_b+\Delta_p}.
\label{eq:F_def}
\end{equation}

对固定的 $p$，先对 $j$ 做收缩：
\begin{equation}
A^p_{ib}
=
\sum_j U^p_{ij}X_{jb}.
\label{eq:A_def}
\end{equation}
把 $X$ 看作形状为 $n_o\times n_v$ 的矩阵，把 $U^p$ 看作形状为 $n_o\times n_o$ 的矩阵，则
\begin{equation}
A^p = U^p X.
\label{eq:A_mat}
\end{equation}

然后乘上能量分母：
\begin{equation}
H^p_{ib}
=
F^p_{ib}A^p_{ib}.
\label{eq:H_def}
\end{equation}
矩阵记号为 Hadamard product：
\begin{equation}
H^p
=
F^p \Had (U^p X).
\label{eq:H_mat}
\end{equation}

最后对 $b$ 做收缩：
\begin{equation}
Y^{(1,p)}_{ia}
=
\sum_b H^p_{ib}V^p_{ab}.
\label{eq:Y1_contract}
\end{equation}
因为 $V^p$ 的矩阵元素是 $V^p_{ab}$，所以上式对应的矩阵乘法为
\begin{equation}
Y^{(1,p)}
=
H^p (V^p)^{\trans}.
\label{eq:Y1_mat}
\end{equation}

因此第一项的完整矩阵表达式为
\begin{equation}
\boxed{
Y^{(1,p)}
=
\left[
F^p\Had (U^pX)
\right]
(V^p)^{\trans}
}.
\label{eq:Y1_final}
\end{equation}

\section{第二项的矩阵化}

第二项为
\begin{equation}
Y^{(2,p)}_{ia}
=
\sum_{jb}
U^p_{ij}V^p_{ab}
\frac{1}{\omega+\eps_j-\eps_a+\Delta_p}
X_{jb}.
\label{eq:second_term}
\end{equation}

定义分母矩阵
\begin{equation}
G^p_{ja}
=
\frac{1}{\omega+\eps_j-\eps_a+\Delta_p}.
\label{eq:G_def}
\end{equation}

注意这里的分母依赖 $(j,a)$，而不是第一项中的 $(i,b)$。
因此第二项更自然的计算顺序是先对 $b$ 做收缩：
\begin{equation}
B^p_{ja}
=
\sum_b X_{jb}V^p_{ab}.
\label{eq:B_def}
\end{equation}

由于 $(V^p)^{\trans}_{ba}=V^p_{ab}$，所以
\begin{equation}
B^p
=
X(V^p)^{\trans}.
\label{eq:B_mat}
\end{equation}

然后乘分母：
\begin{equation}
C^p_{ja}
=
G^p_{ja}B^p_{ja},
\end{equation}
即
\begin{equation}
C^p
=
G^p\Had \left[X(V^p)^{\trans}\right].
\label{eq:C_mat}
\end{equation}

最后对 $j$ 做收缩：
\begin{equation}
Y^{(2,p)}_{ia}
=
\sum_j U^p_{ij}C^p_{ja}.
\label{eq:Y2_contract}
\end{equation}
矩阵形式为
\begin{equation}
Y^{(2,p)}
=
U^p C^p.
\label{eq:Y2_mat}
\end{equation}

所以第二项为
\begin{equation}
\boxed{
Y^{(2,p)}
=
U^p
\left[
G^p\Had \left(X(V^p)^{\trans}\right)
\right]
}.
\label{eq:Y2_final}
\end{equation}

由于原式中第二项带负号，总贡献是
\begin{equation}
\boxed{
Y
=
\sum_p
\left\{
\left[
F^p\Had (U^pX)
\right]
(V^p)^{\trans}
-
U^p
\left[
G^p\Had \left(X(V^p)^{\trans}\right)
\right]
\right\}.
}
\label{eq:main_matrix_formula}
\end{equation}

这就是适合高效实现的核心公式。

\section{RI 张量的矩阵化与 $U^p,V^p$ 的构造}

原始 RI 张量可以在不同二指标空间上展平为矩阵。
定义
\begin{align}
\Roo_{(ij),\mu} &= R_{ij,\mu},
\\
\Rvv_{(ab),\mu} &= R_{ab,\mu},
\\
\Ria_{p,\mu} &= R_{p,\mu}=R_{IA,\mu}.
\end{align}

如果一次处理一个 $p$，则
\begin{align}
U^p_{(ij)}
&=
\sum_{\mu}
\Roo_{(ij),\mu}
\Ria_{p,\mu},
\\
V^p_{(ab)}
&=
\sum_{\mu}
\Rvv_{(ab),\mu}
\Ria_{p,\mu}.
\end{align}

如果一次处理一个 $p$ 的 block，记这个 block 为 $P$，大小为 $B=|P|$。
令
\begin{equation}
\Ria_P \in \mathbb{R}^{B\times n_{\mathrm{aux}}},
\end{equation}
则可以通过两个 GEMM 同时得到所有 $p\in P$ 的 $U^p$ 和 $V^p$：
\begin{align}
U_{\mathrm{flat}}^P
&=
\Roo(\Ria_P)^{\trans},
&
U_{\mathrm{flat}}^P
&\in \mathbb{R}^{n_o^2\times B},
\label{eq:U_gemm}
\\
V_{\mathrm{flat}}^P
&=
\Rvv(\Ria_P)^{\trans},
&
V_{\mathrm{flat}}^P
&\in \mathbb{R}^{n_v^2\times B}.
\label{eq:V_gemm}
\end{align}

然后进行 reshape：
\begin{align}
U_{\mathrm{flat}}^P[(ij),p]
&\longrightarrow
U^p_{ij},
\\
V_{\mathrm{flat}}^P[(ab),p]
&\longrightarrow
V^p_{ab}.
\end{align}

这里的行列排序可以根据实际程序的内存布局自由选择。
只要保证 reshape 后的 $U^p_{ij}$ 与 $V^p_{ab}$ 指标解释一致即可。

\section{直接逐 $p$ 实现的算法}

逐 $p$ 实现最容易检查正确性。
算法如下。

\begin{algorithm}[H]
\caption{逐 $p=(I,A)$ 的 RI-matvec 实现}
\begin{algorithmic}[1]
\Require 输入矩阵 $X\in\mathbb{R}^{n_o\times n_v}$；
         RI 矩阵 $\Roo,\Rvv,\Ria$；
         轨道能量 $\eps_i,\eps_a,\eps_I,\eps_A$；
         频率 $\omega$。
\Ensure 输出矩阵 $Y\in\mathbb{R}^{n_o\times n_v}$。
\State $Y\gets 0$
\For{每一个 $p=(I,A)$}
    \State $\Delta_p\gets \eps_I-\eps_A$
    \State $U^p_{ij}\gets \sum_{\mu}R_{ij,\mu}R_{p,\mu}$
    \State $V^p_{ab}\gets \sum_{\mu}R_{ab,\mu}R_{p,\mu}$
    \State $F^p_{ib}\gets \left(\omega+\eps_i-\eps_b+\Delta_p\right)^{-1}$
    \State $G^p_{ja}\gets \left(\omega+\eps_j-\eps_a+\Delta_p\right)^{-1}$
    \State $A^p\gets U^pX$
    \State $H^p\gets F^p\Had A^p$
    \State $Y\gets Y+H^p(V^p)^{\trans}$
    \State $B^p\gets X(V^p)^{\trans}$
    \State $C^p\gets G^p\Had B^p$
    \State $Y\gets Y-U^pC^p$
\EndFor
\State \Return $Y$
\end{algorithmic}
\end{algorithm}

在代码中，$F^p$ 与 $G^p$ 都是形状为 $n_o\times n_v$ 的矩阵。
它们的数值形式相同：
\begin{equation}
D^p_{\alpha\beta}
=
\frac{1}{\omega+\eps_{\alpha}^{o}-\eps_{\beta}^{v}+\Delta_p}.
\end{equation}
在第一项中这个矩阵解释为 $F^p_{ib}$；
在第二项中解释为 $G^p_{ja}$。
因为 $i,j$ 属于同一个 occupied 空间，$a,b$ 属于同一个 virtual 空间，所以两者可以共享同一个分母矩阵。

\section{按 $p$ 分块的张量化实现}

为了提升效率，实际程序中通常不会逐个 $p$ 计算 $U^p,V^p$，而是按 $p$ 的 block 处理。

对一个 block $P$，大小为 $B$，先通过两个 GEMM 得到：
\begin{align}
U_{pij}, \qquad p\in P,
\\
V_{pab}, \qquad p\in P.
\end{align}

然后第一项可写为
\begin{equation}
A_{pib}
=
\sum_j U_{pij}X_{jb}.
\label{eq:batched_A}
\end{equation}
这是一个 batch 维度为 $B$ 的矩阵乘法：
\begin{equation}
A_p=U_pX.
\end{equation}

定义
\begin{equation}
D_{pib}
=
\frac{1}{\omega+\Delta_p+\eps_i-\eps_b}.
\label{eq:batched_D}
\end{equation}
然后
\begin{equation}
H_{pib}=D_{pib}A_{pib}.
\end{equation}
第一项 block 贡献为
\begin{equation}
Y^{(1,P)}_{ia}
=
\sum_{p\in P}\sum_b H_{pib}V_{pab}.
\label{eq:batched_Y1}
\end{equation}

第二项中：
\begin{equation}
B_{pja}
=
\sum_b X_{jb}V_{pab}.
\label{eq:batched_B}
\end{equation}
然后
\begin{equation}
C_{pja}=D_{pja}B_{pja}.
\end{equation}
第二项 block 贡献为
\begin{equation}
Y^{(2,P)}_{ia}
=
\sum_{p\in P}\sum_j U_{pij}C_{pja}.
\label{eq:batched_Y2}
\end{equation}

因此 block 总贡献为
\begin{equation}
Y
\gets
Y
+
\sum_{p\in P}
\left[
D_p\Had (U_pX)
\right]V_p^{\trans}
-
\sum_{p\in P}
U_p
\left[
D_p\Had (XV_p^{\trans})
\right].
\label{eq:block_update}
\end{equation}

用 Einstein 求和记号可以写为
\begin{align}
A_{pib} &= \mathrm{einsum}(\texttt{"pij,jb->pib"}, U, X),
\\
H_{pib} &= D_{pib}A_{pib},
\\
Y_{ia} &\gets Y_{ia}
+
\mathrm{einsum}(\texttt{"pib,pab->ia"}, H, V),
\\
B_{pja} &= \mathrm{einsum}(\texttt{"jb,pab->pja"}, X,V),
\\
C_{pja} &= D_{pja}B_{pja},
\\
Y_{ia} &\gets Y_{ia}
-
\mathrm{einsum}(\texttt{"pij,pja->ia"}, U,C).
\end{align}

\section{Python 风格伪代码}

下面给出接近实际实现的伪代码。
这里假设：
\begin{itemize}
  \item \texttt{X.shape = (no, nv)}；
  \item \texttt{R\_oo.shape = (no*no, naux)}；
  \item \texttt{R\_vv.shape = (nv*nv, naux)}；
  \item \texttt{R\_IA.shape = (nIA, naux)}；
  \item \texttt{eps\_occ.shape = (no,)}；
  \item \texttt{eps\_vir.shape = (nv,)}；
  \item \texttt{Delta\_IA.shape = (nIA,)}。
\end{itemize}

\begin{verbatim}
Y = zeros((no, nv))

for P in blocks_of_IA:

    # RIA_P: shape (B, naux)
    RIA_P = R_IA[P, :]

    # Two large GEMMs
    # U_flat: shape (no*no, B)
    # V_flat: shape (nv*nv, B)
    U_flat = R_oo @ RIA_P.T
    V_flat = R_vv @ RIA_P.T

    # Reshape to batched matrices
    # U: shape (B, no, no)
    # V: shape (B, nv, nv)
    U = U_flat.T.reshape(B, no, no)
    V = V_flat.T.reshape(B, nv, nv)

    Delta = Delta_IA[P]   # shape (B,)

    # Denominator:
    # D[p,i,a] = 1 / (omega + Delta[p] + eps_occ[i] - eps_vir[a])
    D = 1.0 / (
        omega
        + Delta[:, None, None]
        + eps_occ[None, :, None]
        - eps_vir[None, None, :]
    )

    # First term:
    # A[p,i,b] = sum_j U[p,i,j] X[j,b]
    A = einsum("pij,jb->pib", U, X)

    # H[p,i,b] = D[p,i,b] A[p,i,b]
    H = D * A

    # Y[i,a] += sum_{p,b} H[p,i,b] V[p,a,b]
    Y += einsum("pib,pab->ia", H, V)

    # Second term:
    # Bmat[p,j,a] = sum_b X[j,b] V[p,a,b]
    Bmat = einsum("jb,pab->pja", X, V)

    # C[p,j,a] = D[p,j,a] Bmat[p,j,a]
    C = D * Bmat

    # Y[i,a] -= sum_{p,j} U[p,i,j] C[p,j,a]
    Y -= einsum("pij,pja->ia", U, C)
\end{verbatim}

如果希望更接近 BLAS，可以将三个 \texttt{einsum} 操作分别替换为：
\begin{itemize}
  \item batched GEMM: $A_p=U_pX$；
  \item batched GEMM: $B_p=XV_p^{\trans}$；
  \item batched GEMM 或循环 GEMM 累加：$Y\mathrel{+}=H_pV_p^{\trans}$ 与 $Y\mathrel{-}=U_pC_p$。
\end{itemize}

\section{正确性证明}

下面证明式 \eqref{eq:main_matrix_formula} 与原始求和式 \eqref{eq:original} 完全等价。

\subsection{第一项}

由式 \eqref{eq:Y1_final}，
\begin{equation}
Y^{(1,p)}
=
\left[
F^p\Had (U^pX)
\right]
(V^p)^{\trans}.
\end{equation}
取其 $(i,a)$ 元素：
\begin{align}
Y^{(1,p)}_{ia}
&=
\sum_b
\left[
F^p\Had (U^pX)
\right]_{ib}
(V^p)^{\trans}_{ba}
\\
&=
\sum_b
F^p_{ib}
(U^pX)_{ib}
V^p_{ab}.
\end{align}
展开 $(U^pX)_{ib}$：
\begin{equation}
(U^pX)_{ib}
=
\sum_j U^p_{ij}X_{jb}.
\end{equation}
于是
\begin{align}
Y^{(1,p)}_{ia}
&=
\sum_b
F^p_{ib}
\left(
\sum_j U^p_{ij}X_{jb}
\right)
V^p_{ab}
\\
&=
\sum_{jb}
U^p_{ij}V^p_{ab}
F^p_{ib}
X_{jb}.
\end{align}
代入
\begin{equation}
F^p_{ib}
=
\frac{1}{\omega+\eps_i-\eps_b+\Delta_p},
\end{equation}
得到
\begin{equation}
Y^{(1,p)}_{ia}
=
\sum_{jb}
U^p_{ij}V^p_{ab}
\frac{1}{\omega+\eps_i-\eps_b+\Delta_p}
X_{jb}.
\end{equation}
这正是原始表达式中的第一项。

\subsection{第二项}

由式 \eqref{eq:Y2_final}，
\begin{equation}
Y^{(2,p)}
=
U^p
\left[
G^p\Had \left(X(V^p)^{\trans}\right)
\right].
\end{equation}
取其 $(i,a)$ 元素：
\begin{align}
Y^{(2,p)}_{ia}
&=
\sum_j
U^p_{ij}
\left[
G^p\Had \left(X(V^p)^{\trans}\right)
\right]_{ja}
\\
&=
\sum_j
U^p_{ij}
G^p_{ja}
\left(X(V^p)^{\trans}\right)_{ja}.
\end{align}
展开最后一项：
\begin{equation}
\left(X(V^p)^{\trans}\right)_{ja}
=
\sum_b X_{jb}(V^p)^{\trans}_{ba}
=
\sum_b X_{jb}V^p_{ab}.
\end{equation}
因此
\begin{align}
Y^{(2,p)}_{ia}
&=
\sum_j
U^p_{ij}
G^p_{ja}
\sum_b X_{jb}V^p_{ab}
\\
&=
\sum_{jb}
U^p_{ij}V^p_{ab}
G^p_{ja}
X_{jb}.
\end{align}
代入
\begin{equation}
G^p_{ja}
=
\frac{1}{\omega+\eps_j-\eps_a+\Delta_p},
\end{equation}
得到
\begin{equation}
Y^{(2,p)}_{ia}
=
\sum_{jb}
U^p_{ij}V^p_{ab}
\frac{1}{\omega+\eps_j-\eps_a+\Delta_p}
X_{jb}.
\end{equation}
原始表达式中第二项带负号，所以
\begin{equation}
Y_{ia}
=
\sum_p
\left[
Y^{(1,p)}_{ia}
-
Y^{(2,p)}_{ia}
\right],
\end{equation}
这正好恢复式 \eqref{eq:compact_original}。
再代入 $U^p,V^p$ 的定义式 \eqref{eq:U_def}--\eqref{eq:V_def}，即可恢复原始表达式 \eqref{eq:original}。

\section{计算复杂度分析}

设：
\begin{itemize}
  \item occupied 数为 $n_o$；
  \item virtual 数为 $n_v$；
  \item 辅助基数为 $n_{\mathrm{aux}}$；
  \item 内部 $IA$ pair 数为 $n_{IA}$；
  \item 每个 block 的大小为 $B$。
\end{itemize}

\subsection{构造 $U,V$ 的代价}

对一个 block $P$，构造 $U$ 与 $V$ 的两个 GEMM 代价分别为
\begin{align}
\Roo(\Ria_P)^{\trans}
&:
O(n_o^2 n_{\mathrm{aux}} B),
\\
\Rvv(\Ria_P)^{\trans}
&:
O(n_v^2 n_{\mathrm{aux}} B).
\end{align}

因此每个 block 的 RI 投影代价为
\begin{equation}
O\left[
B n_{\mathrm{aux}}(n_o^2+n_v^2)
\right].
\end{equation}

\subsection{对 $X$ 做变换的代价}

对每个 $p$，核心矩阵乘法包括：
\begin{align}
U^pX &: O(n_o^2n_v),
\\
\left[F^p\Had(U^pX)\right](V^p)^{\trans}
&: O(n_on_v^2),
\\
X(V^p)^{\trans}
&: O(n_on_v^2),
\\
U^p\left[G^p\Had(X(V^p)^{\trans})\right]
&: O(n_o^2n_v).
\end{align}

因此每个 $p$ 的 matvec 代价为
\begin{equation}
O\left(n_o^2n_v+n_on_v^2\right),
\end{equation}
带一个约为 $2$ 的常数因子。

全部 $p$ 的总 matvec 代价为
\begin{equation}
O\left[
n_{IA}(n_o^2n_v+n_on_v^2)
\right].
\end{equation}

\subsection{内存需求}

不显式形成四指标张量时，一个 block 的主要存储为：
\begin{align}
U &: O(Bn_o^2),
\\
V &: O(Bn_v^2),
\\
D &: O(Bn_on_v),
\\
A,H,B,C &: O(Bn_on_v).
\end{align}

所以 block 级内存量级为
\begin{equation}
O\left[
B(n_o^2+n_v^2+n_on_v)
\right].
\end{equation}

相比显式存储四指标对象的
\begin{equation}
O(n_o^2n_v^2),
\end{equation}
该方案通常更加可行，而且可以通过调节 block size $B$ 控制内存。

\section{实现注意事项}

\subsection{指标排序}

实际程序中 $R_{ij,\mu}$、$R_{ab,\mu}$、$R_{IA,\mu}$ 可能已经以不同方式展平。
例如可以使用：
\begin{align}
(ij) &= i n_o+j,
\\
(ab) &= a n_v+b,
\\
(IA) &= I n_v+A.
\end{align}
也可以使用 Fortran-order 的排列方式：
\begin{align}
(ij) &= j n_o+i,
\\
(ab) &= b n_v+a.
\end{align}
数学推导不依赖具体排序。
但实现中必须保证：
\begin{equation}
U_{\mathrm{flat}}[(ij),p]\rightarrow U^p_{ij},
\qquad
V_{\mathrm{flat}}[(ab),p]\rightarrow V^p_{ab}
\end{equation}
的 reshape 与后续 GEMM 的矩阵指标解释一致。

\subsection{能量分母的广播}

对每个 block，可以使用广播构造
\begin{equation}
D_{p i a}
=
\frac{1}{
\omega+\Delta_p+\eps_i-\eps_a
}.
\end{equation}
在代码中对应：
\begin{verbatim}
D = 1.0 / (
    omega
    + Delta[:, None, None]
    + eps_occ[None, :, None]
    - eps_vir[None, None, :]
)
\end{verbatim}

该矩阵在第一项中作为 $D_{pib}$ 使用，在第二项中作为 $D_{pja}$ 使用。
因为第二维始终是 occupied 指标，第三维始终是 virtual 指标，所以可以共用同一个数组。

\subsection{复频率与阻尼}

如果 $\omega$ 含有虚部，例如
\begin{equation}
\omega \rightarrow \omega+i\eta,
\end{equation}
则 $D,F,G$ 变为复数矩阵。
上述公式完全不变，只需要把相关数组和 GEMM 类型改为复数。
这在 damped response 或共振附近计算中尤其自然。

\subsection{是否需要形成 $D$}

若 $B n_o n_v$ 较大，可以不显式保存整个 $D$。
例如在逐 $p$ 循环中临时生成 $D^p$；
或者在 block 内按 occupied/virtual 子块进一步切分。
这可以把内存从
\begin{equation}
O(Bn_on_v)
\end{equation}
进一步降低到子块规模。

\subsection{GEMM 与 batched GEMM}

推荐的高性能实现策略是：
\begin{enumerate}
  \item 使用大 GEMM 构造 $U_{\mathrm{flat}}^P$ 与 $V_{\mathrm{flat}}^P$；
  \item 使用 batched GEMM 计算 $U_pX$；
  \item 使用 batched GEMM 计算 $XV_p^{\trans}$；
  \item 使用 batched GEMM 或循环 GEMM 累加 $H_pV_p^{\trans}$ 与 $U_pC_p$；
  \item 使用 elementwise kernel 完成分母缩放。
\end{enumerate}

如果 $n_o,n_v$ 较小而 $B$ 较大，batched GEMM 通常比逐个调用 GEMM 更有效。
如果 $n_o,n_v$ 较大，逐个 GEMM 的效率也可能已经足够高。

\section{最终公式总结}

整个 matvec 可以总结为以下步骤。

首先，对每个 $p=(I,A)$ 定义：
\begin{align}
\Delta_p &= \eps_I-\eps_A,
\\
U^p_{ij} &= \sum_{\mu}R_{ij,\mu}R_{p,\mu},
\\
V^p_{ab} &= \sum_{\mu}R_{ab,\mu}R_{p,\mu},
\\
D^p_{\alpha\beta}
&=
\frac{1}{\omega+\Delta_p+\eps_\alpha^{o}-\eps_\beta^{v}}.
\end{align}

然后输出为
\begin{equation}
\boxed{
Y
=
\sum_p
\left\{
\left[
D^p\Had(U^pX)
\right](V^p)^{\trans}
-
U^p
\left[
D^p\Had(X(V^p)^{\trans})
\right]
\right\}.
}
\end{equation}

其中第一处 $D^p$ 的指标解释为 $D^p_{ib}$，第二处 $D^p$ 的指标解释为 $D^p_{ja}$。
二者数值上是同一个 occupied--virtual 能量分母矩阵。

这个表达式避免了四指标张量 $\widetilde W_{ij,ab}$ 的显式构造，将原始求和转化为：
\begin{itemize}
  \item 两次 RI 投影 GEMM；
  \item 若干 batched GEMM；
  \item 能量分母的逐元素缩放；
  \item 对 $IA$ block 的累加。
\end{itemize}

因此它适合在现有 RI 框架、BLAS/LAPACK、GPU batched GEMM 或自定义张量 contraction kernel 中实现。

\end{document}
